\chapter{Projects and production planning}
\label{chap:planificacion}

\section{CPM: precedence and the critical path}
CPM operates on an acyclic activity network: an activity starts only after
all its predecessors have finished. It does not allocate limited resources.
The forward pass calculates early starts and finishes; the backward pass
calculates late starts and finishes. Total float is late start minus early
start. A critical activity has zero float.

In \path{manual/examples/plan_cpm.json}, A takes 3 days, B takes 2, and both
precede C, which takes 4. Each activity has a normal cost of 10, a crash
duration one unit shorter, and a crash cost of 15.
\begin{lstlisting}[language=bash]
swift run qsb validate-project-json manual/examples/plan_cpm.json
swift run qsb solve-project-json manual/examples/plan_cpm.json --backend native
swift run qsb solve-cpm-crash-json manual/examples/plan_cpm.json 6 --backend native
\end{lstlisting}
\begin{figure}[htbp]
\centering
\begin{tikzpicture}
\node[qsb activity,line width=1.4pt] (A) at (0,1) {A\\3};
\node[qsb activity] (B) at (0,-1) {B\\2};
\node[qsb activity,line width=1.4pt] (C) at (4,0) {C\\4};
\draw[qsb selected] (A) -- (C);
\draw[qsb alternative] (B) -- (C);
\end{tikzpicture}

\caption{CPM activity network: the number below each name is its duration in days. Thick strokes identify critical activities and precedence; the dashed stroke shows the other precedence.}
\label{fig:cpm}
\end{figure}

The normal project finishes in 7 days. A and C are critical; B has float 1.
The total normal cost is 30. Crashing to a deadline of 6 reduces C from 4
to 3, with incremental cost 5 and total cost 35. Reducing A could also
achieve the same cost; do not assume that the plan is unique.

Crash optimization assumes linear costs between normal and minimum
durations, with slope $(C_{crash}-C_{normal})/(t_{normal}-t_{crash})$ when
reduction is possible. It uses continuous times and an educational LP that
considers paths; the number of paths can grow rapidly. It is not a policy
for hiring resources and does not include indirect project costs.
Deadlines shorter than what is achievable must be treated as infeasible,
never as instructions to violate precedence constraints.

\section{PERT: expected duration and uncertainty}
PERT uses optimistic $a$, most likely $m$, and pessimistic $b$ estimates,
ordered $a\leq m\leq b$. The classical approximation calculates
\[
 E[T]=(a+4m+b)/6,\qquad \operatorname{Var}(T)=((b-a)/6)^2.
\]
In \path{plan_pert.json}, the triples are A(2,3,4), B(1,2,3), and C(3,4,5),
with the same precedence relationships as in the CPM example.
\begin{lstlisting}[language=bash]
swift run qsb solve-project-json manual/examples/plan_pert.json --backend native
\end{lstlisting}
\begin{figure}[htbp]
\centering
\begin{tikzpicture}
\node[qsb activity,line width=1.4pt] (A) at (0,1) {A\\(2,3,4)\\$E[T]=3$};
\node[qsb activity] (B) at (0,-1) {B\\(1,2,3)\\$E[T]=2$};
\node[qsb activity,line width=1.4pt] (C) at (4,0) {C\\(3,4,5)\\$E[T]=4$};
\draw[qsb selected] (A) -- (C);
\draw[qsb alternative] (B) -- (C);
\end{tikzpicture}

\caption{PERT network: each activity shows its optimistic, most likely, and pessimistic triple and its expected duration in days. Thick strokes identify the critical path calculated using those means.}
\label{fig:pert}
\end{figure}

Scheduling with expected durations again gives a duration of 7, with A and
C critical. The reported variance is $2/9$ and the standard deviation is
$0.471405$ days. In this example, there is a single critical path, and
summing its variances assumes independence. The implementation sums the
variances of all activities marked critical; if there are multiple critical
paths, that total is neither the variance of a single path nor that of the
maximum over them. The critical path may change between realizations;
replacing random durations with their means does not calculate the expected
maximum exactly. The \texttt{exact} label describes the calculation using
those expected durations, not an exact probabilistic prediction of the
completion date. QSB does not return a probability of meeting a deadline
or simulate that uncertainty in this module.

\subsection{CPM and PERT in macOS}
Import the project envelope or create a CPM/PERT definition. The activity
editor collects names, predecessors, and durations; CPM adds costs, and
PERT adds the three estimates. Missing predecessors and cycles appear in
Validation. Switching methods preserves data for both methods after
initialization: check the estimates when returning to PERT.
Copy Table/Paste Table lets you transfer activities from a spreadsheet.
Run produces duration, critical activities, early/late bars, float values,
and a timing table. Review JSON for the variance and exact values.
The CLI crash operation is a separate workflow; the normal view does not
imply that a deadline has been optimized. Native and Validate are
available; External does not currently solve project models.

\section{Aggregate planning}
Aggregate planning balances demand across periods through regular
production, overtime, subcontracting, inventory, and backorders.
When capacity depends on workforce, it adds hiring and layoffs.
A basic balance is
\[
 I_t-B_t=I_{t-1}-B_{t-1}+R_t+O_t+S_t-D_t.
\]
The objective sums production, idle capacity, workforce, inventory, and
backorder costs according to the data. QSB normalizes the historical simple,
transportation, and LP methods into a single continuous LP formulation;
the \texttt{method} field preserves the input identity rather than selecting
three different native algorithms.

The example \path{plan_aggregate.json} has two months, demand 4 and 6,
regular capacity 5 each month, regular cost 2 per unit, overtime capacity
up to 2 at cost 5, inventory cost 1, and no backorders or subcontracting.
It does not depend on workforce and starts with no inventory.
\begin{lstlisting}[language=bash]
swift run qsb validate-aggregate-json manual/examples/plan_aggregate.json
swift run qsb solve-aggregate-json manual/examples/plan_aggregate.json --backend native
\end{lstlisting}
The plan produces 5 in each month; it holds 1 at the end of the first and
zero at the end of the second. It costs $10\cdot2+1=21$. Producing exactly
the first month's demand and using one overtime unit in the second would
cost $9\cdot2+5=23$.
Units must be consistent: capacity per worker is not total capacity, and
\texttt{capacityRequirementPerUnit} converts production into capacity.
Vectors must have the same number of periods; \texttt{null} in optional
maximum fields means no upper bound. Workforce and quantities are
continuous, even when the business needs whole people. Do not round a
solution without rechecking balances and capacity.

In macOS, the period editor separates demand, costs, inventory/sources,
and workforce. Copy Table transfers all 20 period columns; Paste Table
validates the range before replacing it. Run shows stacked production and
demand, inventory, backorders, and a plan table. A backorder appears below
zero in the chart, but its model quantity is nonnegative. Native solves
the educational LP; External uses HiGHS for the same formulation if
installed. Validate does not produce a plan. The optimal objective refers
to the costs and bounds entered, not to omitted criteria such as workforce
stability and its social implications.

\section{MRP: requirements explosion}
MRP starts from a master production schedule (MPS), a bill of materials
(BOM), inventory, scheduled receipts, and lead times. It processes parents
before components: a parent's order release generates gross requirements
for its components, multiplied by the quantity used per unit
\cite{silver1998inventory}.
For an item with safety stock $SS$, gross requirement $G_t$, and scheduled
receipt $SR_t$, net requirement is
$\max(0,G_t+SS-I_{t-1}-SR_t)$. A planned receipt covers that requirement,
and its release is shifted backward by the lead time.

\path{plan_mrp.json} has weeks S0--S3. Five units of P must be received in
S3; P uses two units of C per unit. Both have a one-week lead time, zero
initial inventory and safety stock, the LFL rule, and unlimited capacity.
\begin{lstlisting}[language=bash]
swift run qsb validate-mrp-json manual/examples/plan_mrp.json
swift run qsb solve-mrp-json manual/examples/plan_mrp.json --backend native
\end{lstlisting}
P receives 5 in S3 and releases 5 in S2. C needs 10 in S2, receives 10
then, and releases 10 in S1. Projected inventories remain zero. C's
requirements come from P's release, not from the final demand date.
The BOM must be acyclic, and all its identifiers must exist.

The available rules are LFL, EOQ, LUC, LTC, and PPB. LFL receives the net
requirement; EOQ starts from $\sqrt{2DS/H}$ and receives at least the
current requirement. In EOQ, \texttt{lotSizingLookahead} is interpreted as
a lot multiple \(m=\max(1,\text{configured value})\), not as a number of
periods: the quantity is
\(\max(\text{requirement},m\lceil EOQ/m\rceil)\).
If annual holding cost is zero, it receives the net requirement.
LUC compares unit cost; LTC and PPB currently use the same selector, which
minimizes the absolute difference between holding cost and ordering cost
within the configured horizon. These are forward-looking rules, not a
global optimization of multilevel costs. QSB labels the MRP document
\texttt{heuristic}, even when the particular case uses LFL.
Finite capacity generates \texttt{capacityExcess}; it does not level or
reschedule orders. If a release falls before the horizon, QSB accumulates
it in the first bucket using $\max(0,t-L)$; that record requires operational
review and does not prove that the lead time can be met.

In macOS, Item Master collects parameters, Schedules collects quantities
by period, and Bill of Materials collects parent/component relationships.
Select an item to edit or read its schedule; long horizons are grouped for
navigation. Copy/Paste Table always exchanges the complete range for its
section. The solution shows requirements, inventory, receipts, releases,
and capacity excesses by item, with JSON as an audit record. Native and
Validate are available; External is not implemented for MRP.

\section{Exercises}
\begin{enumerate}
\item In CPM, increase B from 2 to 4. Recalculate the critical path and float.
\item In PERT, increase only A's pessimistic estimate. Compare duration and
      variance; explain which assumptions limit a probabilistic interpretation.
\item Raise aggregate inventory cost from 1 to 4. Compare early production
      with overtime.
\item Add initial inventory of 4 units of C in MRP and check its net
      requirement. Then limit C's capacity in S2 to 5 and explain the excess.
\end{enumerate}
